% ch52_presentable.tex — Volume IV Chapter 16
% Retrofitted with full Vol I-III pedagogy: cycle stages, etymology, dialogs.

\chapter{Presentable $\infty$-Categories}

\begin{overviewbox}
\term{Presentable $\infty$-categories} are the working class of
``manageable'' $\infty$-categories: those that are cocomplete and
generated under small colimits by a small ``presentation.'' Examples:
$\mathcal{S}$ (spaces), $\mathcal{P}(\mathcal{C})$ (presheaves on a
small $\mathcal{C}$), and most $\infty$-categories arising in derived
algebraic geometry, homotopy theory, and topos theory.

This chapter develops the theory:
\begin{itemize}
  \item \term{$\kappa$-filtered colimits} for a regular cardinal
        $\kappa$ — the building blocks of accessibility.
  \item \term{Accessible $\infty$-categories} — those generated under
        $\kappa$-filtered colimits by a small set of $\kappa$-compact
        objects.
  \item \term{Presentable $\infty$-categories} — accessible $+$
        cocomplete. Equivalently, accessible localizations of $\mathcal{P}(\mathcal{C})$.
\end{itemize}

The capstone is the \term{$\infty$-Adjoint Functor Theorem} for
presentable $\infty$-categories: a functor $R : \mathcal{D} \to
\mathcal{C}$ admits a left adjoint iff $R$ preserves all small
$\infty$-limits and is \emph{accessible}. This is the practical engine
for proving existence of $\infty$-categorical free constructions, base
changes, sheafifications, and the like.

A reader who completes this chapter has acquired the standard ``working
infrastructure'' assumed throughout Lurie's \emph{Higher Topos Theory}
(Chapters 5–6).
\end{overviewbox}


% ═══════════════════════════════════════════════════════════════════════════
\section{Accessibility}
\label{sec:accessibility}
% ═══════════════════════════════════════════════════════════════════════════

\subsection{Formalizing}

\begin{nameorigin}
\term{Accessible} and \term{presentable} categories were introduced by
\emph{Jiří Adámek} and \emph{Jiří Rosický} in their 1994 book
\emph{Locally Presentable and Accessible Categories}. The terminology:
\emph{accessible} means ``reachable from a small set of generators
under $\kappa$-filtered colimits'' (the category is ``accessible''
from a small generating set); \emph{presentable} (or \emph{locally
presentable}) means accessible + cocomplete (i.e., ``every object has
a small presentation in terms of generators and relations''). The
$\infty$-version was developed by Lurie in HTT Ch.~5 with the same
philosophy.
\end{nameorigin}

\begin{definition}[$\kappa$-filtered]
\label{def:kappa-filtered}
For a regular cardinal $\kappa$, a quasi-category $\mathcal{K}$ is
\term{$\kappa$-filtered} if every diagram $K \to \mathcal{K}$ of size
$< \kappa$ admits a cocone in $\mathcal{K}$.

A $\omega$-filtered ($\kappa = \aleph_0$) quasi-category is just a
filtered quasi-category in the sense of Chapter~49.
\end{definition}

\begin{nameorigin}
\term{$\kappa$-compact} (also called \term{$\kappa$-presentable}) is from
topology / set theory: a topological space is \emph{compact} if every
open cover has a finite subcover; an object of a category is \emph{$\kappa$-compact}
if every map from it into a $\kappa$-filtered colimit factors through a
stage. The terminology was lifted into category theory by Gabriel and
Ulmer (1971). The classical $\omega$-compact (= compact) objects in
$\mathbf{Top}$ are the finite CW-complexes — generalizing the
topological compactness exactly.
\end{nameorigin}

\begin{definition}[$\kappa$-compact object]
\label{def:kappa-compact}
An object $x$ of a cocomplete $\infty$-category $\mathcal{C}$ is
\term{$\kappa$-compact} (or \term{$\kappa$-presentable}) if
$\mathrm{Map}_\mathcal{C}(x, -) : \mathcal{C} \to \mathcal{S}$ preserves
$\kappa$-filtered colimits.

The full subcategory of $\kappa$-compact objects is $\mathcal{C}^\kappa
\subseteq \mathcal{C}$.
\end{definition}

\begin{example_thm}[Compactness in $\mathcal{S}$]
\label{ex:compact-in-S}
An object of $\mathcal{S}$ is $\omega$-compact iff it is equivalent to a
finite CW complex (a finitely-presented Kan complex). For larger
$\kappa$, $\kappa$-compact spaces are exactly Kan complexes with
$\kappa$-bounded simplices.
\end{example_thm}

\begin{definition}[Accessible $\infty$-category]
\label{def:accessible}
A quasi-category $\mathcal{C}$ is \term{accessible} if there is a regular
cardinal $\kappa$ and a small set $\{x_i\}_{i \in I}$ of $\kappa$-compact
objects such that every object of $\mathcal{C}$ is a $\kappa$-filtered
colimit of the $x_i$.

Equivalently: $\mathcal{C}$ has $\kappa$-filtered colimits and the
inclusion $\mathcal{C}^\kappa \hookrightarrow \mathcal{C}$ extends along
$\kappa$-filtered colimits to an equivalence
$\mathrm{Ind}^\kappa(\mathcal{C}^\kappa) \simeq \mathcal{C}$.
\end{definition}

\begin{howtosay}
$\mathrm{Ind}^\kappa(\mathcal{D})$ \sayit{``the $\kappa$-Ind completion of
$\mathcal{D}$''}: freely add $\kappa$-filtered colimits to $\mathcal{D}$.
\end{howtosay}

\begin{theorem}[Accessibility-via-Ind]
\label{thm:ind-accessibility}
For a small $\infty$-category $\mathcal{D}$, the universal $\kappa$-filtered
colimit completion $\mathrm{Ind}^\kappa(\mathcal{D})$ exists, is the
$\kappa$-filtered-colimit closure of $\mathcal{D}$ in $\mathcal{P}(\mathcal{D})$,
and is universal among accessible $\infty$-categories admitting $\mathcal{D}$
as their $\kappa$-compact subcategory.
\end{theorem}


% ═══════════════════════════════════════════════════════════════════════════
\section{Presentable $\infty$-Categories}
\label{sec:presentable}
% ═══════════════════════════════════════════════════════════════════════════

\subsection{Formalizing}

\begin{definition}[Presentable $\infty$-category]
\label{def:presentable}
A quasi-category $\mathcal{C}$ is \term{presentable} (or \term{locally
presentable}) if it is:
\begin{enumerate}
  \item Cocomplete (admits all small $\infty$-colimits).
  \item Accessible (a $\kappa$-Ind completion of a small $\infty$-category, for some $\kappa$).
\end{enumerate}
\end{definition}

\begin{theorem}[Characterizations of presentability]
\label{thm:presentable-characterizations}
A quasi-category $\mathcal{C}$ is presentable iff any (equivalently, all) of:
\begin{enumerate}
  \item It is an \term{accessible localization} of some
        $\mathcal{P}(\mathcal{D})$ — i.e., a reflective subcategory
        whose reflector preserves $\kappa$-filtered colimits for some
        $\kappa$.
  \item It is cocomplete and there is a small set of objects
        $\{x_i\}$ that generates $\mathcal{C}$ under small colimits.
  \item There exists a small $\mathcal{D}$ and a small set of morphisms
        $W$ in $\mathcal{P}(\mathcal{D})$ such that $\mathcal{C} \simeq
        \mathcal{P}(\mathcal{D})[W^{-1}]$ — a Bousfield-localization
        presentation.
\end{enumerate}
\end{theorem}

\begin{remark}[The class $\mathrm{Pr}^L$]
The category $\mathrm{Pr}^L$ has presentable $\infty$-categories as
objects and \emph{colimit-preserving} (left adjoint) functors as
morphisms. Lurie's HTT establishes that $\mathrm{Pr}^L$ is itself an
$\infty$-category, in fact a closed symmetric monoidal $\infty$-category
under the \emph{Lurie tensor product}. This last is the working substrate
of $\infty$-operads (Chapter~55) and stable homotopy theory.
\end{remark}

\subsection{Reorganizing: foundational examples}

\begin{example_thm}[$\mathcal{S}$ is presentable]
\label{ex:S-presentable}
$\mathcal{S}$ is presentable: $\omega$-compact objects are finite CW
complexes, and every space is an $\omega$-filtered colimit of finite ones.
In fact, $\mathcal{S} \simeq \mathrm{Ind}^\omega(\mathcal{S}^\omega)$.
\end{example_thm}

\begin{example_thm}[$\mathcal{P}(\mathcal{C})$ is presentable]
\label{ex:PC-presentable}
For any small $\mathcal{C}$, $\mathcal{P}(\mathcal{C})$ is presentable.
$\omega$-compact objects: presheaves with finite support (representable
or finite colimits of representables). Generation: representables.
\end{example_thm}

\begin{example_thm}[Modules over a ring]
\label{ex:modules}
$\mathrm{Mod}_R$ for a connective $\mathbb{E}_\infty$-ring spectrum $R$
is presentable. $\omega$-compact: finitely generated projective
$R$-modules. Generation: $R$ itself. This is the principal example of
practical higher algebra.
\end{example_thm}


% ═══════════════════════════════════════════════════════════════════════════
\section{The $\infty$-Adjoint Functor Theorem}
\label{sec:aft}
% ═══════════════════════════════════════════════════════════════════════════

\subsection{Formalizing}

\begin{theorem}[$\infty$-AFT, presentable case]
\label{thm:aft-full}
Let $\mathcal{C}, \mathcal{D}$ be presentable $\infty$-categories. A
functor $R : \mathcal{D} \to \mathcal{C}$ admits a left adjoint iff:
\begin{enumerate}
  \item $R$ preserves all small $\infty$-limits.
  \item $R$ is accessible: preserves $\kappa$-filtered colimits for some
        regular cardinal $\kappa$.
\end{enumerate}
Dually, $L : \mathcal{C} \to \mathcal{D}$ admits a right adjoint iff $L$
preserves all small $\infty$-colimits.
\end{theorem}

\begin{proof}[Proof strategy]
The proof is HTT §5.5.2 (Theorem 5.5.2.9). The strategy:

\emph{Existence of left adjoint.} For each $y \in \mathcal{D}$, the
functor $\mathrm{Map}_\mathcal{C}(-, R y) \circ R : \mathcal{D} \to \mathcal{S}$
preserves limits and is accessible. By the \term{Adjoint Functor Theorem
for $\mathcal{S}$-valued functors}, it is representable; the representing
object is the value of $L$ at the corresponding object of $\mathcal{C}$.

\emph{Uniqueness} follows from $\infty$-Yoneda.

\emph{Necessity of conditions} is Theorem~50.\ref{thm:rapl-lapc-infty}
(RAPL) plus the observation that left adjoints to accessible functors
are accessible.

The dual statement uses the symmetric structure: right adjoints exist
iff left adjoints to opposite functors do. \qed
\end{proof}

\begin{example_thm}[Building free objects]
\label{ex:free-construction}
$\mathrm{Mod}_R \to \mathcal{S}$, sending a module to its underlying space,
preserves all limits and is accessible. The adjoint functor theorem
produces a left adjoint $\mathcal{S} \to \mathrm{Mod}_R$: the \emph{free
$R$-module} functor, $S \mapsto R \otimes \Sigma^\infty_+ S$.
\end{example_thm}


% ═══════════════════════════════════════════════════════════════════════════
\section{Bousfield Localization and Reflective Subcategories}
\label{sec:bousfield}
% ═══════════════════════════════════════════════════════════════════════════

\subsection{Formalizing}

\begin{nameorigin}
\term{Bousfield localization} is named for \emph{Aldridge Knight ``Pete''
Bousfield} (1941--2020), American mathematician who pioneered the
construction in his 1975 paper ``The Localization of Spectra with
Respect to Homology.'' Bousfield's insight was that one can localize an
entire homotopy category at a chosen class of maps without leaving the
ambient category — by working with the reflective subcategory of
``$W$-local objects.'' The construction is the homotopy-theoretic
generalization of the Gabriel-Zisman localization (Ch.~40), and it
extends naturally to the $\infty$-categorical setting via reflective
subcategories of presentable $\infty$-cats.
\end{nameorigin}

\begin{definition}[Reflective subcategory]
\label{def:reflective}
A full subcategory $\mathcal{D} \hookrightarrow \mathcal{C}$ is
\term{reflective} if the inclusion has a left adjoint $L : \mathcal{C}
\to \mathcal{D}$. The composition $\mathcal{C} \xrightarrow{L} \mathcal{D}
\hookrightarrow \mathcal{C}$ is then an idempotent monad (a \term{localization}).
\end{definition}

\begin{theorem}[Bousfield localization]
\label{thm:bousfield}
For any presentable $\mathcal{C}$ and small set $W$ of morphisms of
$\mathcal{C}$, there is a reflective subcategory $L_W \mathcal{C}
\hookrightarrow \mathcal{C}$ whose objects are the \term{$W$-local} ones:
$x$ such that $\mathrm{Map}_\mathcal{C}(w, x)$ is an equivalence for
every $w \in W$.

The reflection $L_W : \mathcal{C} \to L_W \mathcal{C}$ inverts every
morphism in $W$ and is universal among colimit-preserving functors out
of $\mathcal{C}$ that do.
\end{theorem}

\begin{remark}[Bousfield localization and presentability]
Every presentable $\infty$-category arises (up to equivalence) as a
Bousfield localization of some $\mathcal{P}(\mathcal{C})$: this is
Theorem~\ref{thm:presentable-characterizations}'s characterization (3).
In particular, the $\infty$-category of sheaves on a Grothendieck site is
the Bousfield localization of presheaves of spaces at the covers.
\end{remark}


% ═══════════════════════════════════════════════════════════════════════════
\section{Exercises}
\label{sec:presentable-exercises}
% ═══════════════════════════════════════════════════════════════════════════

\begin{exercisesbox}
\begin{qa}
\qaitem{Define $\kappa$-filtered.}{A quasi-category $\mathcal{K}$ is $\kappa$-filtered if every diagram of size $< \kappa$ has a cocone.}
\qaitem{What's a $\kappa$-compact object?}{$x \in \mathcal{C}$ is $\kappa$-compact if $\mathrm{Map}_\mathcal{C}(x, -)$ preserves $\kappa$-filtered colimits.}
\qaitem{Define accessible.}{$\mathcal{C}$ is accessible if it equals $\mathrm{Ind}^\kappa(\mathcal{C}^\kappa)$ for some $\kappa$ and the small subcategory $\mathcal{C}^\kappa$.}
\qaitem{Define presentable.}{Cocomplete and accessible. Equivalent: accessible localization of $\mathcal{P}(\mathcal{C})$ for some small $\mathcal{C}$.}
\qaitem{What are $\omega$-compact objects in $\mathcal{S}$?}{Finite CW complexes (equivalently, finitely generated Kan complexes).}
\qaitem{Why is $\mathcal{P}(\mathcal{C})$ presentable?}{Cocomplete by construction (limits and colimits computed pointwise); accessible because representables are compact and generate under colimits.}
\qaitem{What is $\mathrm{Pr}^L$?}{The $\infty$-category of presentable $\infty$-categories with colimit-preserving (left adjoint) functors as morphisms. Closed symmetric monoidal under Lurie tensor product.}
\qaitem{State the $\infty$-adjoint functor theorem.}{For presentable $\mathcal{C}, \mathcal{D}$: $R : \mathcal{D} \to \mathcal{C}$ has a left adjoint iff $R$ preserves $\infty$-limits and is accessible. Dually for left adjoints.}
\qaitem{Why does the AFT work?}{Map-Yoneda: limit-preserving accessible $\mathcal{D} \to \mathcal{S}$ is representable; this gives the adjoint value-by-value.}
\qaitem{What is Bousfield localization?}{For presentable $\mathcal{C}$ and small $W \subseteq \mathrm{Mor}(\mathcal{C})$, the reflective subcategory $L_W \mathcal{C} \subseteq \mathcal{C}$ on $W$-local objects.}
\qaitem{Define $W$-local.}{$x$ is $W$-local iff $\mathrm{Map}(w, x)$ is an equivalence of spaces for every $w \in W$.}
\qaitem{Why is every sheaf $\infty$-category a Bousfield localization?}{Sheaves are presheaves local with respect to covers: $\mathrm{Sh}(X) = L_{\mathrm{covers}} \mathcal{P}(\mathrm{Open}(X))$.}
\qaitem{What is $\mathrm{Ind}^\kappa(\mathcal{D})$?}{The free $\kappa$-filtered colimit completion of $\mathcal{D}$: representables plus all $\kappa$-filtered colimits of them.}
\qaitem{Why do compactness and accessibility differ from $1$-categories?}{They don't, conceptually — the definitions are direct analogues — but in the $\infty$-setting ``size'' and ``finite generation'' are subtler because of the higher-dimensional structure of mapping spaces.}
\qaitem{What's a reflective subcategory?}{Full subcategory $\mathcal{D} \hookrightarrow \mathcal{C}$ with a left adjoint $L : \mathcal{C} \to \mathcal{D}$. The composition $\mathcal{C} \to \mathcal{D} \to \mathcal{C}$ is an idempotent endofunctor.}
\qaitem{Foundational example of localization?}{The localization of $\mathbf{sSet}$ at weak equivalences gives $\mathcal{S}$. The localization of $\mathcal{P}(\mathbf{Aff})$ at étale covers gives the $\infty$-topos of étale stacks.}
\qaitem{How does presentability help with adjoint existence?}{Presentability ensures that AFT-style ``representability'' arguments terminate (the small-presentation provides a stopping point for transfinite induction).}
\qaitem{Are all left adjoints in $\mathrm{Pr}^L$ accessible?}{Yes — accessibility of a functor reflects in/out of accessible $\infty$-cats, so left adjoints between presentables are automatically accessible.}
\qaitem{What's the difference between presentable and locally presentable?}{The terms are interchangeable in most modern literature; ``locally presentable'' emphasizes the underlying $1$-categorical analogue (Adámek–Rosický).}
\qaitem{What's a generator?}{A small set $\{x_i\}$ of objects such that every object of $\mathcal{C}$ is a colimit of $x_i$'s. Compact objects of an accessible $\infty$-cat form a generator.}
\qaitem{Why is the Lurie tensor product important?}{It makes $\mathrm{Pr}^L$ symmetric monoidal, allowing tensor products of presentable $\infty$-categories — central to symmetric monoidal $\infty$-categories (Ch 55) and ring spectra.}
\qaitem{What's the next chapter?}{Ch 53: Stable $\infty$-Categories. Presentability + stability gives ``presentable stable $\infty$-categories'' = $\mathrm{Pr}^L_{\mathrm{Stab}}$ — the working category for stable homotopy theory and derived algebraic geometry.}
\qaitem{Why does presentability matter in practice?}{Because nearly every $\infty$-category arising in modern mathematics — sheaves, modules, schemes, types — is presentable. The AFT is then a free tool.}
\end{qa}
\end{exercisesbox}


% ═══════════════════════════════════════════════════════════════════════════
\section*{Chapter Summary}
\addcontentsline{toc}{section}{Chapter Summary}
% ═══════════════════════════════════════════════════════════════════════════

\begin{summarybox}
\textbf{$\kappa$-filtered and $\kappa$-compact.}\quad $\kappa$-filtered:
diagrams of size $< \kappa$ have cocones. $\kappa$-compact: $\mathrm{Map}(x, -)$
preserves $\kappa$-filtered colimits.

\textbf{Accessible.}\quad $\mathcal{C} \simeq \mathrm{Ind}^\kappa(\mathcal{C}^\kappa)$
for some $\kappa$ and small $\mathcal{C}^\kappa$. Generated under
$\kappa$-filtered colimits by a small set.

\textbf{Presentable.}\quad Cocomplete + accessible. Equivalent
characterizations: accessible localization of $\mathcal{P}(\mathcal{C})$;
small generator under colimits; Bousfield localization $\mathcal{P}(\mathcal{C})[W^{-1}]$.

\textbf{$\infty$-AFT.}\quad For presentable $\mathcal{C}, \mathcal{D}$:
$R$ has a left adjoint iff $R$ preserves $\infty$-limits and is accessible.
Engine of free constructions throughout higher algebra.

\textbf{Bousfield localization.}\quad For any small $W \subseteq
\mathrm{Mor}(\mathcal{C})$ in a presentable $\mathcal{C}$, the
$W$-local subcategory $L_W \mathcal{C}$ is reflective with reflector
inverting $W$. Every presentable $\infty$-cat is a Bousfield localization
of some $\mathcal{P}(\mathcal{C})$.

\textbf{$\mathrm{Pr}^L$.}\quad The $\infty$-category of presentable
$\infty$-categories with colimit-preserving functors. Closed symmetric
monoidal under Lurie tensor product.
\end{summarybox}


% ═══════════════════════════════════════════════════════════════════════════
\section*{Formal Restatement}
\addcontentsline{toc}{section}{Formal Restatement}
% ═══════════════════════════════════════════════════════════════════════════

\begin{formalrestatement}
\textbf{Accessibility.}\quad $\mathcal{C}$ accessible $\iff
\mathcal{C} \simeq \mathrm{Ind}^\kappa(\mathcal{D})$ for some small
$\mathcal{D}$ and regular $\kappa$.

\medskip
\textbf{Presentability.}\quad
$\mathcal{C}$ presentable $\iff$ cocomplete and accessible.

\medskip
\textbf{Characterizations (equivalent).}\quad
\begin{itemize}
  \item Accessible localization of $\mathcal{P}(\mathcal{D})$.
  \item Cocomplete with a small generating set.
  \item $\mathcal{P}(\mathcal{D})[W^{-1}]$ for small $\mathcal{D}, W$.
\end{itemize}

\medskip
\textbf{$\infty$-AFT.}\quad For presentable $\mathcal{C}, \mathcal{D}$:
\begin{equation*}
  R : \mathcal{D} \to \mathcal{C} \text{ has left adjoint} \iff
  R \text{ preserves } \infty\text{-limits and is accessible}.
\end{equation*}
Dually, $L$ has right adjoint $\iff$ $L$ preserves $\infty$-colimits.

\medskip
\textbf{Bousfield localization.}\quad $\mathcal{C}$ presentable,
$W$ small $\Rightarrow L_W \mathcal{C} \hookrightarrow \mathcal{C}$
reflective; $L_W$ universal among $W$-inverting colimit-preserving
functors out of $\mathcal{C}$.

\medskip
\textbf{$\mathrm{Pr}^L$.}\quad The $\infty$-category of presentable
$\infty$-cats with colimit-preserving functors. Closed symmetric monoidal
under Lurie tensor.
\end{formalrestatement}
